\section{Generalisation Experiments}
\label{sec:generalisation}

\subsection{Zero-Shot Transfer to Unicorn}

Unicorn~\cite{wang2025unicorn} is applied to the same 8iVFB sequences at
five rate points (R2--R6), covering a wider bitrate range than PCGCv2
including extreme compression (R5--R6 at $\bpp < 0.025$).
The post-processor is applied without any retraining or fine-tuning.

Table~\ref{tab:unicorn-results} reports the results.

\begin{table}[t]
  \centering
  \caption{%
    Zero-shot cross-codec transfer to Unicorn on 8iVFB.
    All values averaged over 4 sequences.
    The model was trained exclusively on PCGCv2-decoded data.
  }
  \label{tab:unicorn-results}
  \begin{tabular}{lccrr}
    \toprule
    Rate & Approx. bpp & Edge $\delta$ (dB) & Flat $\delta$ (dB) \\
    \midrule
    R2 & 0.17 & $-0.27$ & $+0.34$ \\
    R3 & 0.071 & $-0.72$ & $-0.02$ \\
    R4 & 0.052 & $-0.01$ & $+0.30$ \\
    R5 & 0.023 & $\mathbf{+1.09}$ & $+0.40$ \\
    R6 & 0.017 & $\mathbf{+1.36}$ & --- \\
    \bottomrule
  \end{tabular}
\end{table}

\begin{figure}[t]
  \centering
  \includegraphics[width=0.72\textwidth]{unicorn_transfer.png}
  \caption{%
    Zero-shot transfer results to Unicorn.
    Green bars indicate positive PSNR improvement; red bars indicate degradation.
    Transfer succeeds at extreme rates (R5--R6, $\bpp < 0.025$)
    where oversmoothing is codec-agnostic, but degrades at moderate rates
    (R2--R3) where Unicorn's distortion pattern differs from PCGCv2's.
  }
  \label{fig:unicorn-transfer}
\end{figure}

The results show a clear bitrate threshold.
At moderate rates (R2--R4), Unicorn's distortion profile differs from PCGCv2's:
the model applies corrections calibrated for PCGCv2-style oversmoothing to a
different geometry distribution, causing slight degradation at edges.
At extreme rates (R5--R6), the oversmoothing is severe enough to be
codec-agnostic --- both PCGCv2 and Unicorn produce heavily blurred surfaces
at $\bpp < 0.025$ --- and the model's corrections transfer.

This finding has a practical implication: for scenarios where extreme geometric
fidelity is required at very low bitrates (e.g., constrained streaming),
a codec-agnostic post-processor is a viable solution.

\subsection{Zero-Shot Transfer to MVUB}
\label{ssec:mvub}

The MVUB dataset~\cite{loop2016mvub} provides upper-body captures at
resolution $G = 512$ (two sequences: \textit{david} and \textit{ricardo}),
representing an out-of-distribution test for both codec and dataset.
The model is applied without retraining.

Table~\ref{tab:mvub-results} reports the results, using BPP-corrected
rate assignment (the per-frame BPP from the actual Unicorn compressed
MVUB bitstreams).

\begin{table}[t]
  \centering
  \caption{%
    Zero-shot transfer to MVUB dataset (PCGCv2 decoded, $G=512$).
    Results averaged over 10 frames per sequence.
    BPP-corrected rate used for FiLM conditioning.
  }
  \label{tab:mvub-results}
  \begin{tabular}{llrrr}
    \toprule
    Sequence & Rate & Edge $\delta$ (dB) & Flat $\delta$ (dB) & D1 $\delta$ (dB) \\
    \midrule
    \multirow{3}{*}{david}
      & r2 & $\mathbf{+0.92}$ & $\mathbf{+0.47}$ & $-0.24$ \\
      & r4 & $-0.46$ & $-0.05$ & $-0.42$ \\
      & r6 & $-0.09$ & $-0.09$ & $-0.10$ \\
    \midrule
    \multirow{3}{*}{ricardo}
      & r2 & $\mathbf{+1.03}$ & $\mathbf{+0.43}$ & $-0.88$ \\
      & r4 & $-0.50$ & $-0.14$ & $-0.48$ \\
      & r6 & $-0.12$ & $-0.12$ & $-0.13$ \\
    \bottomrule
  \end{tabular}
\end{table}

The pattern mirrors the Unicorn transfer: positive gains at r2 (low rate),
negative at r4/r6 (moderate-high rates).
At r2, both \textit{david} ($+0.92$~dB) and \textit{ricardo} ($+1.03$~dB)
show substantial edge improvement at the rate where oversmoothing is most
severe, suggesting the correction holds across dataset and resolution.
The negative D1 $\delta$ at r2 is notable: the model improves edge quality
but slightly degrades the aggregate D1, reflecting a geometric trade-off
between edge recovery (adding sub-voxel displacements to edge points) and
flat-region precision.
This reveals a limitation of aggregate D1 as a metric for evaluating
post-processing that specifically targets non-uniform artifacts.

% ============================================================
